\paragraph{C.4aa. A countermodel to an unrestricted graph-transfer principle (Proved).}
Let $\Lambda=\mathbb Z_{(2)}$, let $U=\mathbb Z^2$, and consider
\[
 q_1=x^2,\qquad q_2=y^2,\qquad q_3=(2x+y)^2.
\]
There is no finite three-coloured graph $G$ and no positive integer
$d$ admitting a common $\Lambda$-orthogonal direct-sum decomposition
\[
 H_1(G,\Lambda)=\iota(U_\Lambda)\oplus\ker\beta,
 \qquad \beta\iota=1,
\]
for which the restriction of the colour-$s$ graph form is $dq_s$ for
all three colours. Nevertheless, these three forms are the monodromy
forms of an algebraic family of principally polarized abelian surfaces,
every geometric fibre of which has an algebraic minimal curve class.
Thus a graph-transfer conclusion of this form cannot follow from
algebraicity alone for arbitrary integral monodromy cones.

\begin{proof}
Put
\[
 L=\begin{pmatrix}1&0\\0&1\\2&1\end{pmatrix},
 \qquad R=\begin{pmatrix}-2&-1&1\end{pmatrix}.
\]
The sequence $0\to\mathbb Z^2\xrightarrow L\mathbb Z^3
\xrightarrow R\mathbb Z\to0$ is split exact, and every row $L_s$
maps $U$ onto $\mathbb Z$. Suppose a graph splitting exists.
Positivity makes the projected coedge of an edge of colour $s$
equal to $v_eL_s$, with $v_e\in\Lambda$. The common orthogonal
retraction gives the profile identity
\[
 \Phi(h)=dL\beta(h)\quad(h\in H_1(G,\Lambda)),
 \qquad \sum_{e\text{ of colour }s}v_e^2=d.
\]
Here $\Phi$ puts $v_e$ in coordinate $s$. Indeed, pairing $h$
with $\iota(u)$ in the colour-$s$ form and using orthogonality
gives $L_su$ times the $s$th profile coordinate on one side, and
$d(L_su)(L_s\beta h)$ on the other; surjectivity of $L_s$ gives
the asserted identity.

Write $d=2^j d'$ with $d'$ odd and put $a=j+2$. The coordinate
rounding argument proved in C.4z applies with three scalar blocks
in place of its three two-dimensional blocks. We specify all its
inputs here. Replace an edge of colour $s$ by a positive integral
axis path whose length is congruent to $v_e$ modulo $2^a$, retaining
the original scalar flux $v_e$ on each step. The period identity
maps the resulting graph into the coordinate-axis graph of
\[
 \mathbb R^3/P,
 \qquad P=2^jL\mathbb Z^2+2^{j+2}\mathbb Z^3.
\]
Its coordinate-$s$ profile $M_s\in\Lambda$ satisfies
$M_s\equiv d\pmod{2^a}$. The coordinate rounding map that
stretches the first unit interval of each block of length $2^j$
and contracts the remaining intervals is $P$-equivariant and
homotopic to the identity. It preserves each $M_s$ exactly:
choose $u$ with $L_su=1$ and read the $s$th displacement of the
associated closed embedded chain. The other colours never move
in coordinate $s$. Since $P\subset2^j\mathbb Z^3$, interiors
of distinct coarse axis edges do not meet. At an internal
subdivision vertex, closure for every $u$ forces the scalar
flux on the adjacent steps to coincide, again because $L_s$
is onto. Collapsing the $2^j$ steps and rescaling divides all
three profiles by $2^j$.

Consequently there is an axis graph with vertices
\[
 \mathbb Z^3/(L\mathbb Z^2+4\mathbb Z^3)
                  \simeq\mathbb Z/4\mathbb Z
\]
carrying three scalar-flux chains with profiles
$M'_s\equiv d'\pmod4$. Their combination with coefficients
$L_su$ is closed for every $u$. The three positive axis steps
have displacements $-2,-1,1$ in this quotient.
Multiply the first coordinate of the flow-boundary relation
at a vertex $t$ by $f(t)=t^2\pmod4$ and sum over all vertices.
Colour one contributes zero because
$f(t-2)-f(t)\equiv0\pmod4$. Colour two has zero first flow
coordinate. A colour-three edge of flux $v$ contributes
\[
 2v\bigl(f(t+1)-f(t)\bigr)
       =2v(2t+1)\equiv2v\pmod4.
\]
The boundary sum is therefore $2M'_3\equiv2d'\pmod4$,
whereas closure makes it zero. This contradiction proves
the assertion for every positive level $d$. This argument,
as in C.4z, uses a genuine common retraction and not merely
an embedding with the prescribed restrictions.

For the geometric countermodel, apply
\citet[Proposition~2.29 and Corollary~2.30]{frEGFSMumfordSurvey2026}
to the three positive-semidefinite integral matrices
\[
 \begin{pmatrix}1&0\\0&0\end{pmatrix},\quad
 \begin{pmatrix}0&0\\0&1\end{pmatrix},\quad
 \begin{pmatrix}4&2\\2&1\end{pmatrix}.
\]
The cited result supplies an algebraic projective family of
principally polarized abelian surfaces over the complement
of an snc divisor in a smooth three-dimensional base, with
these exact coordinate monodromies. On any principally
polarized abelian surface, the minimal curve class is the
class of a theta divisor itself. Hence it is an effective
integral algebraic curve class on the geometric generic
fibre as well. No assertion about smoothness or irreducibility
of that theta divisor is needed for algebraicity. The
algebraic minimal class coexists with the nonexistence of
the asserted common graph at every level, proving the
claimed limitation.
\end{proof}

The same example has pairwise distinct one-dimensional monodromy
images, hence pairwise zero intersection. Thus the observation that
a node involving two original base colours has zero projected
coedge does not, by itself, produce a stable-curve graph over the
original base. A special geometric argument for the correlated
rank-two cone in C.4z would require additional information.
